\section{The \texttt{Slit} McXtrace Component}
Rectangular/circular slit

\subsection*{Identification}
\begin{itemize}
  \item \textbf{Author:} Erik Knudsen
  \item \textbf{Origin:} DTU Physics
  \item \textbf{Date:} June 16, 2009
\end{itemize}

\subsection*{Description}
Based on Slit-comp by Kim Lefmann and Henrik Roennow A simple rectangular or circular slit. You may either specify the radius (circular shape), which takes precedence, or rectangular bounds. No transmission around the slit is allowed.

Example: Slit(xmin=-0.01, xmax=0.01, ymin=-0.01, ymax=0.01) Slit(radius=0.01)

\subsection*{Input parameters}
Parameters in \textbf{boldface} are required; the others are optional.

\begin{longtable}{p{0.22\textwidth}p{0.12\textwidth}p{0.46\textwidth}p{0.14\textwidth}}
\toprule
\textbf{Name} & \textbf{Unit} & \textbf{Description} & \textbf{Default} \\
\midrule
\endhead
xmin & m & Lower x bound. & 0 \\
xmax & m & Upper x bound. & 0 \\
ymin & m & Lower y bound. & 0 \\
ymax & m & Upper y bound. & 0 \\
radius & m & Radius of slit in the z=0 plane, centered at Origin. & 0 \\
xwidth & m & Width of slit. Overrides xmin,xmax. & 0.001 \\
yheight & m & Height of slit. Overrides ymin,ymax. & 0.001 \\
dist & m & Distance from slit plane to plane containing resampling target. & 0 \\
focus\_xw & m & Width of resampling window. & 0 \\
focus\_yh & m & Height of resampling window. & 0 \\
focus\_x0 & m & Centre (x) of resampling window. & 0 \\
focus\_y0 & m & Centre (y) of resampling window. & 0 \\
\bottomrule
\end{longtable}

\subsection*{Links}
\begin{itemize}
  \item Component source code found in file \texttt{Slit.comp}.
\end{itemize}
\IfFileExists{optics/Slit_static.tex}{\input{optics/Slit_static.tex}}{}